\chapter{Existence Theory for the Ricci Flow}

\section*{5.1 \quad Ricci flow is not parabolic}

Consider
\[
\frac{\partial g}{\partial t} = Q(g) := -2\Ric(g),
\tag{5.1.1}
\]
on the bundle $E=\Sym^2 T^*\M$. The linearisation of \textup{(5.1.1)} is
\[
\frac{\partial h}{\partial t}=Lh:= [DQ(g)]h = \Delta_{\mathcal L}h+\mathcal L_{(\dG(h))^\#}g,
\]
so we compute its principal symbol. Since $\Delta_{\mathcal L}=\Delta$ up to lower-order terms,
\[
\sigma(\Delta_{\mathcal L})(x,\xi)h = |\xi|^2 h.
\]
Moreover, using $\dG(h)=\delta h+\frac12 d(\tr h)$ and the identity
\[
\mathcal L_{\omega^\#}g(X,W)=\nabla\omega(X,W)+\nabla\omega(W,X),
\]
we find
\[
\sigma\!\Bigl(h\mapsto \mathcal L_{(\dG(h))^\#}g\Bigr)(x,\xi)h
=-\xi\otimes h(\xi^\sharp,\cdot)-h(\xi^\sharp,\cdot)\otimes\xi+(\xi\otimes\xi)\tr h.
\]
Hence
\[
\sigma(L)(x,\xi)h
=|\xi|^2h-\xi\otimes h(\xi^\sharp,\cdot)-h(\xi^\sharp,\cdot)\otimes\xi+(\xi\otimes\xi)\tr h.
\]
Setting $h=\xi\otimes\xi$ gives $\sigma(L)(x,\xi)h=0$, so the linearised Ricci flow is not parabolic.

\begin{remark}
\label{topping-ricci-flow-linearisation-not-parabolic}
\dcref{5.1.1}
\group{topping-ricci-flow-existence}
\source{topping:page-0060,topping:page-0061}
\uses{topping:variation-ricci,topping:lie-derivative-symmetric-gradient}
The Ricci-flow linearisation has principal symbol
\[
\sigma(L)(x,\xi)h
=|\xi|^2h-\xi\otimes h(\xi^\sharp,\cdot)-h(\xi^\sharp,\cdot)\otimes\xi+(\xi\otimes\xi)\tr h,
\]
and vanishes on $h=\xi\otimes\xi$. In particular, the Ricci flow equation is not parabolic.
\end{remark}

\section*{5.2 \quad Short-time existence and uniqueness: the DeTurck trick}

Let $T\in\Gamma(\Sym^2T^*\M)$ be fixed, smooth, and positive definite, and let $P$ be defined by
\[
P(g)=-2\Ric(g)+\mathcal L_{T^{-1}\dG(T)}^{\#}g.
\tag{5.2.1}
\]
Linearising the correction term shows that the second-order part of $DP(g)$ is just $\Delta$; hence $P$ is parabolic at every metric.

\begin{theorem}[Short time existence]
\label{topping-ricci-flow-short-time-existence}
\dcref{5.2.1}
\group{topping-ricci-flow-existence}
\source{topping:page-0062,topping:page-0063}
\uses{topping:variation-ricci,topping:lie-derivative-symmetric-gradient}
Given a smooth metric $g_0$ on a closed manifold $\M$, there exist $\epsilon>0$ and a smooth family of metrics $g(t)$ for $t\in[0,\epsilon]$ such that
\[
\begin{cases}
\dfrac{\partial g}{\partial t}=-2\Ric(g) & t\in[0,\epsilon],\\
g(0)=g_0.
\end{cases}
\tag{5.2.2}
\]
\end{theorem}

\begin{proof}
The DeTurck-modified equation $\partial_t g=P(g)$ is parabolic, so the short-time existence theory for quasilinear parabolic systems on closed manifolds gives a solution with initial data $g_0$. Integrating the time-dependent vector field
\[
X=-(T^{-1}\dG(T))^\#
\]
produces diffeomorphisms $\psi_t$ with $\psi_0=\id$. Pulling back the DeTurck solution by $\psi_t$ cancels the Lie-derivative term and yields a Ricci flow with the same initial metric.
\end{proof}

\begin{theorem}[Uniqueness of solutions (forwards and backwards)]
\label{topping-ricci-flow-uniqueness}
\dcref{5.2.2}
\group{topping-ricci-flow-existence}
\source{topping:page-0064}
Suppose $g_1(t)$ and $g_2(t)$ are two Ricci flows on a closed manifold $\M$, for $t\in[0,\epsilon]$, $\epsilon>0$. If $g_1(s)=g_2(s)$ for some $s\in[0,\epsilon]$, then $g_1(t)=g_2(t)$ for all $t\in[0,\epsilon]$.
\end{theorem}

\begin{proof}
For each flow, solve the harmonic map flow into a fixed background metric $T$ with initial map $\id$. The corresponding pullbacks of $g_1(t)$ and $g_2(t)$ satisfy the same DeTurck equation with the same initial metric, so uniqueness for the parabolic system forces them to agree. The harmonic-map gauges therefore agree, and pushing back gives $g_1(t)=g_2(t)$ for all $t\in[0,\epsilon]$.
\end{proof}

Combining short-time existence and uniqueness, one may speak of the Ricci flow with initial metric $g_0$ on a maximal time interval $[0,T]$.

\section*{5.3 \quad Curvature blow-up at finite-time singularities}

\begin{lemma}[Metric equivalence]
\label{topping-ricci-flow-metric-equivalence}
\dcref{5.3.2}
\group{topping-ricci-flow-existence}
\source{topping:page-0065}
If $g(t)$ is a Ricci flow for $t\in[0,s]$ and $|\Ric|\le M$ on $\M\times[0,s]$, then
\[
e^{-2Mt}g(0)\le g(t)\le e^{2Mt}g(0)
\]
for all $t\in[0,s]$.
\end{lemma}

\begin{proof}
For any nonzero tangent vector $X$,
\[
\frac{\partial}{\partial t}g(X,X)=-2\Ric(X,X),
\]
so $|\partial_t g(X,X)|\le 2M\,g(X,X)$. Dividing by $g(X,X)$ gives
\[
\left|\frac{\partial}{\partial t}\ln g(X,X)\right|\le 2M.
\]
Integrating in time yields
\[
\left|\ln\frac{g(t)(X,X)}{g(0)(X,X)}\right|\le 2Mt,
\]
which is equivalent to the stated two-sided bound.
\end{proof}

\begin{theorem}[Curvature blows up at a singularity]
\label{topping-ricci-flow-curvature-blows-up}
\dcref{5.3.1}
\group{topping-ricci-flow-existence}
\source{topping:page-0064,topping:page-0065,topping:page-0066,topping:page-0067}
\uses{topping-maxpr-curvature-norm-bound,topping-maxpr-higher-derivatives,topping-ricci-flow-metric-equivalence}
If $\M$ is closed and $g(t)$ is a Ricci flow on a maximal time interval $[0,T)$ with $T<\infty$, then
\[
\sup_{\M}|\Rm|(\cdot,t)\to\infty
\qquad (t\uparrow T).
\]
\end{theorem}

\begin{proof}
Assume instead that $|\Rm|$ stays bounded on $[0,T)$. By Theorem~3.2.11 there is a uniform bound $|\Rm|\le M$ on $\M\times[0,T)$. Since $|\Ric|$ is controlled by $|\Rm|$, the metric equivalence lemma shows that $g(t)$ converges continuously to a positive-definite limit tensor $g(T)$ as $t\uparrow T$.

The derivative estimates from Corollary~3.3.2 bound $\partial_t^\ell \nabla^k\Rm$ on a time slab $[\frac12T,T)$; combining these with the Ricci-flow equation and the commutator formula for $\partial_t$ and $\nabla$ bounds all mixed space-time derivatives of the coordinate coefficients $g_{ij}$ on $[\frac12T,T)$. Thus $g(t)$ extends smoothly to $t=T$.

Now regard $g(T)$ as initial data in the short-time existence theorem. This extends the flow past $T$, contradicting maximality. Therefore $|\Rm|$ must blow up as $t\uparrow T$.
\end{proof}
